\documentclass[12pt,a4paper]{article}

\usepackage[utf8]{inputenc}
\usepackage[T1]{fontenc}
\usepackage[spanish]{babel}
\usepackage{geometry}
\usepackage{array}
\usepackage{tabularx}
\usepackage{longtable}
\usepackage{booktabs}
\usepackage{enumitem}
\usepackage{xcolor}
\usepackage{hyperref}
\usepackage{graphicx}
\usepackage{setspace}
\usepackage{titlesec}
\usepackage{fancyhdr}

\geometry{margin=2.3cm}
\setstretch{1.12}
\setlength{\parskip}{0.45em}
\setlength{\parindent}{0pt}

\definecolor{utpblue}{HTML}{0A2A66}
\definecolor{utpcyan}{HTML}{57E8FF}
\definecolor{softgray}{HTML}{F5F7FB}

\hypersetup{
    colorlinks=true,
    linkcolor=utpblue,
    urlcolor=utpblue
}

\titleformat{\section}{\large\bfseries\color{utpblue}}{\thesection.}{0.5em}{}
\titleformat{\subsection}{\normalsize\bfseries\color{utpblue}}{\thesubsection.}{0.5em}{}

\pagestyle{fancy}
\fancyhf{}
\fancyhead[L]{Proyecto final de Programación 2}
\fancyhead[R]{Robot Explota Globos UTP}
\fancyfoot[C]{\thepage}

% Datos editables
\newcommand{\tituloProyecto}{Plataforma Web para la Gestión del Torneo Robot Explota Globos}
\newcommand{\nombreAplicacion}{Robot Explota Globos UTP}
\newcommand{\integranteUnoNombre}{[Nombre integrante 1]}
\newcommand{\integranteUnoCodigo}{[Código integrante 1]}
\newcommand{\integranteUnoCorreo}{[Correo integrante 1]}
\newcommand{\integranteDosNombre}{[Nombre integrante 2 / opcional]}
\newcommand{\integranteDosCodigo}{[Código integrante 2 / opcional]}
\newcommand{\integranteDosCorreo}{[Correo integrante 2 / opcional]}

\begin{document}

\begin{titlepage}
    \centering
    {\Large Universidad Tecnológica de Pereira\par}
    \vspace{0.4cm}
    {\large Programación 2\par}
    \vspace{1.8cm}
    {\Huge\bfseries \tituloProyecto \par}
    \vspace{0.7cm}
    {\Large Propuesta de proyecto final\par}
    \vspace{1.5cm}
    \begin{tabular}{>{\bfseries}p{4.8cm}p{9.5cm}}
        Nombre de la aplicación: & \nombreAplicacion \\
        Tipo de solución: & Plataforma web administrativa y pública para la operación de un torneo académico de robótica \\
        Curso: & Programación 2 \\
        Fecha: & Mayo de 2026 \\
    \end{tabular}
    \vfill
    {\small
    Nota de adaptación: aunque la rúbrica original referencia aplicaciones construidas con Flet, esta propuesta se basa al 100\% en el proyecto real ya desarrollado y validado por el docente como entrega equivalente. La solución propuesta utiliza Django, PostgreSQL, HTML, CSS y JavaScript como base tecnológica principal.
    \par}
\end{titlepage}

\section{Título del proyecto y nombre de la aplicación}

\textbf{Título del proyecto:} \tituloProyecto

\textbf{Nombre de la aplicación:} \nombreAplicacion

La aplicación corresponde a una plataforma digital para gestionar la inscripción, validación, seguimiento y desarrollo competitivo del torneo \emph{Explota Globos}, evento académico impulsado desde el programa de Ingeniería Electrónica de la Universidad Tecnológica de Pereira.

\section{Información general de los integrantes}

\begin{longtable}{>{\bfseries}p{4.2cm}p{3.5cm}p{6.2cm}}
\toprule
Nombre & Código & Correo \\
\midrule
\endfirsthead
\toprule
Nombre & Código & Correo \\
\midrule
\endhead
\integranteUnoNombre & \integranteUnoCodigo & \integranteUnoCorreo \\
\integranteDosNombre & \integranteDosCodigo & \integranteDosCorreo \\
\bottomrule
\end{longtable}

\textit{Observación:} los campos anteriores quedan listos para ser reemplazados con los datos definitivos antes de la entrega.

\section{Descripción general del proyecto}

El proyecto consiste en una plataforma web para la gestión integral del torneo de robots ``Explota Globos'' de la Universidad Tecnológica de Pereira. El sistema resuelve una necesidad operativa real: centralizar en una sola herramienta el reglamento oficial, el registro de participantes, la validación de asistencia, la organización de categorías por colegios y universidades, y el control administrativo del cuadro del torneo.

La problemática principal que se busca solucionar es la dependencia de procesos manuales o semimanuales en un evento con muchos participantes, lo cual puede generar errores en la inscripción, dificultad para controlar fases, pérdida de trazabilidad y reorganizaciones improvisadas durante la competencia. La plataforma propone una solución robusta, editable y trazable que permita a los organizadores actuar en tiempo real sin perder el historial de lo ocurrido en cada etapa.

El contexto de uso es académico y logístico. La plataforma será utilizada durante la etapa previa al torneo para registrar robots, confirmar asistencia y difundir información oficial, y durante el evento para administrar grupos, repechajes, cruces y resultados desde un panel de control interno.

\section{Objetivo del proyecto}

Desarrollar una plataforma web segura, flexible y visualmente clara que permita administrar de manera centralizada el torneo Robot Explota Globos, automatizando el proceso de inscripción y la generación inicial del formato competitivo, pero conservando al mismo tiempo control manual total para que los organizadores puedan ajustar grupos, fases, estados y resultados según las necesidades reales del evento.

\section{Funciones principales}

\subsection{Módulo público}

\begin{itemize}[leftmargin=1.2cm]
    \item Página principal del torneo con identidad visual del evento.
    \item Publicación del reglamento oficial en una vista embebida.
    \item Sección de patrocinadores con interacción visual y enlaces externos.
    \item Registro diferenciado para colegios y universidades.
    \item Validaciones para evitar duplicidad de robots, correos o documentos.
    \item Flujo de confirmación por correo para validar asistencia antes del torneo.
\end{itemize}

\subsection{Módulo administrativo}

\begin{itemize}[leftmargin=1.2cm]
    \item Panel interno por categorías: colegios y universidades.
    \item Generación automática de estructura competitiva según el número de participantes.
    \item Modo editable para personalizar grupos, clasificados, repechajes y tercer lugar.
    \item Navegación independiente por fases: grupos, repechajes, rondas eliminatorias, semifinal y final.
    \item Modal de edición por participante para cambiar grupo, fase, estado o reintegrarlo manualmente.
    \item Trazabilidad histórica por participante para registrar cambios y resultados.
\end{itemize}

\subsection{Almacenamiento de la información}

La solución plantea un esquema de almacenamiento híbrido, alineado con el proyecto ya construido:

\begin{itemize}[leftmargin=1.2cm]
    \item \textbf{Base de datos principal PostgreSQL:} almacena usuarios, ediciones del torneo, reglas, inscripciones, participantes, equipos, grupos, fases, batallas, estados actuales e historial de cambios.
    \item \textbf{Archivos CSV:} se utilizarán para exportar reportes operativos como listados de inscritos, confirmaciones de asistencia, distribución por categorías y consolidado de resultados.
    \item \textbf{Archivos JSON:} se utilizarán para respaldar configuraciones variables del torneo, estructuras de fases, metadatos de combates, snapshots operativos y trazas de auditoría externa si se requiere intercambio con otros sistemas.
\end{itemize}

De esta manera, se cumple con un manejo de datos que no depende únicamente del registro de inicio de sesión y que además permite flexibilidad para análisis, respaldo y visualización.

\subsection{Operaciones con los datos almacenados}

Sobre la información almacenada, la aplicación realizará operaciones como:

\begin{itemize}[leftmargin=1.2cm]
    \item Filtrado de inscripciones por categoría, estado y edición activa.
    \item Validación de unicidad de robots y participantes.
    \item Confirmación de asistencia y habilitación de participantes para entrar al cuadro oficial.
    \item Generación aleatoria y balanceada de grupos.
    \item Cálculo de clasificados y envío de equipos a repechajes o fases posteriores.
    \item Reasignación manual de participantes en caso de abandono, daño del robot o reorganización del torneo.
    \item Consulta histórica de movimientos y resultados por participante.
\end{itemize}

\subsection{Procesamiento visual de datos}

La aplicación incluirá al menos dos gráficas para apoyar decisiones de organización:

\begin{enumerate}[leftmargin=1.2cm]
    \item \textbf{Gráfica de barras de inscripciones por categoría y estado:} permitirá ver cuántos registros hay en colegios y universidades, diferenciando entre enviados, aprobados, confirmados y pendientes.
    \item \textbf{Gráfica circular o de dona del avance del torneo:} mostrará la distribución de participantes activos, clasificados, eliminados y en repechaje por categoría.
\end{enumerate}

Estas gráficas servirán para análisis administrativo previo y seguimiento durante el evento.

\section{Librerías a utilizar}

Aunque la rúbrica base menciona Flet y Pandas, en esta versión aprobada del proyecto la interfaz principal ya fue construida como aplicación web. Las librerías y herramientas más relevantes son:

\begin{itemize}[leftmargin=1.2cm]
    \item \textbf{Django:} framework principal para el backend, rutas, autenticación, plantillas y manejo del panel administrativo.
    \item \textbf{Pandas:} procesamiento tabular para reportes, análisis de inscripciones, confirmaciones y resultados.
    \item \textbf{psycopg:} conexión entre Django y PostgreSQL.
    \item \textbf{python-dotenv:} carga de variables de entorno y manejo seguro de credenciales.
    \item \textbf{Matplotlib o Plotly:} generación de gráficas estadísticas dentro del módulo administrativo o para exportación de reportes.
\end{itemize}

\section{Audiencia objetivo / usuarios esperados}

La aplicación está dirigida a varios perfiles:

\begin{itemize}[leftmargin=1.2cm]
    \item \textbf{Organizadores del torneo:} requieren control operativo de grupos, fases, resultados y confirmaciones.
    \item \textbf{Docentes y equipo logístico:} necesitan supervisar la información oficial y resolver situaciones excepcionales.
    \item \textbf{Participantes de colegios y universidades:} utilizan la parte pública para registrarse, leer reglas y confirmar asistencia.
    \item \textbf{Aliados institucionales y patrocinadores:} aparecen visibilizados en la plataforma como parte del ecosistema del evento.
\end{itemize}

El perfil esperado combina usuarios administrativos con conocimientos básicos de navegación web y participantes externos que necesitan una experiencia simple, clara y rápida.

\section{Estructura preliminar del proyecto}

En el contexto de una aplicación web, la noción de ``ventanas'' se adapta a \textbf{vistas o páginas}. La estructura preliminar actual del sistema contempla:

\subsection{Cantidad de vistas principales}

Se plantean aproximadamente nueve vistas principales, algunas públicas y otras privadas:

\begin{enumerate}[leftmargin=1.2cm]
    \item Inicio
    \item Reglas
    \item Patrocinadores
    \item Selección de tipo de registro
    \item Formulario de colegios
    \item Formulario de universidades
    \item Confirmación / asistencia
    \item Dashboard administrativo general
    \item Vistas administrativas por fase
\end{enumerate}

\subsection{Función de cada vista}

\begin{longtable}{>{\bfseries}p{4.2cm}p{9.5cm}}
\toprule
Vista & Función principal \\
\midrule
\endfirsthead
\toprule
Vista & Función principal \\
\midrule
\endhead
Inicio & Presentar el evento, dirigir al reglamento, al registro y destacar patrocinadores. \\
Reglas & Mostrar el reglamento oficial embebido dentro de la plataforma. \\
Patrocinadores & Exhibir aliados del evento con logos interactivos y enlaces externos. \\
Selección de registro & Permitir que el usuario elija si se inscribe como colegio o universidad. \\
Formulario de colegios & Capturar los datos del robot y sus integrantes para la categoría colegios. \\
Formulario de universidades & Capturar los datos del robot y sus integrantes para la categoría universidades. \\
Confirmación / asistencia & Gestionar mensajes de recepción del registro y confirmación de asistencia mediante correo. \\
Dashboard administrativo & Mostrar métricas globales, acceso a categorías y control general del torneo. \\
Vistas por fase & Operar grupos, repechajes, llaves eliminatorias y edición manual por participante. \\
\bottomrule
\end{longtable}

\subsection{Colores o temática visual}

La propuesta visual ya construida toma como base una estética tecnológica y electrónica:

\begin{itemize}[leftmargin=1.2cm]
    \item fondo oscuro con contraste alto;
    \item acentos en cian, azul eléctrico y magenta;
    \item tipografía con estética tecnológica;
    \item efectos de hover y animaciones suaves para mejorar percepción de dinamismo;
    \item identidad coherente con Ingeniería Electrónica y con el torneo de robots.
\end{itemize}

\subsection{Componentes gráficos previstos}

En las diferentes vistas se emplearán componentes como:

\begin{itemize}[leftmargin=1.2cm]
    \item barras de navegación responsivas;
    \item tarjetas informativas;
    \item formularios con validación;
    \item modales de edición rápida;
    \item grids adaptativos;
    \item visualizadores PDF;
    \item galerías y secciones interactivas de patrocinadores;
    \item tablas y tarjetas de estado en el panel administrativo;
    \item gráficas estadísticas para reportes de inscripciones y avance del torneo.
\end{itemize}

\section{Conclusión}

La propuesta presentada no es una idea hipotética desligada de la práctica, sino la formalización académica de un sistema que ya cuenta con una base funcional importante. El proyecto combina una interfaz pública para participantes con un panel administrativo avanzado para gestionar un torneo real, incluyendo inscripción, confirmación, organización de fases, edición manual y trazabilidad histórica.

Desde la perspectiva de Programación 2, el proyecto demuestra integración de estructuras de datos, validaciones, persistencia, interfaces gráficas, procesamiento de información y diseño de flujos de usuario orientados a resolver un problema concreto. Por tanto, representa una propuesta sólida, pertinente y escalable para ser presentada como proyecto final del curso.

\end{document}
